% year: תשפ"ה
% type: בוחן
% semester: א
% source: exams/בחנים/תשפ״ה, בוחן, 9141.pdf
% part: ב
% number: 5
% points: 10
% categories: functions, multiple-choice

\begin{question}
תהי $f(x)=3^{-3x+3}$. אילו מהטענות הבאות נכונה?
\begin{enumerate}
  \item הפונקציה אינה מונוטונית.
  \item הפונקציה מונוטונית יורדת ממש.
  \item הפונקציה מונוטונית עולה ממש.
  \item הפונקציה מונוטונית אך לא ממש.
\end{enumerate}
\end{question}

\begin{solution}
\textbf{התשובה הנכונה: (ב)}.

$f$ היא הרכבה $f=h\circ u$ של $u(x)=-3x+3$ ו-$h(t)=3^t$. הפונקציה $u$ יורדת ממש (פונקציה לינארית עם שיפוע $-3<0$) והפונקציה $h$ עולה ממש (פונקציה מעריכית עם בסיס $3>1$).

לכן אם $x_1<x_2$ אז $u(x_1)>u(x_2)$, ומכאן $3^{u(x_1)}>3^{u(x_2)}$, כלומר $f(x_1)>f(x_2)$: $f$ יורדת ממש על $\R$.
(לחלופין: $f(x)=27\cdot\left(\frac{1}{27}\right)^{x}$, פונקציה מעריכית עם בסיס $\frac1{27}<1$ כפול קבוע חיובי.)

\textbf{למה האחרות שגויות:} (א) — הראינו שהיא מונוטונית. (ג) — למשל $f(0)=27>f(1)=1$, ולכן אינה עולה. (ד) — "מונוטונית אך לא ממש" אומר שיש $x_1<x_2$ עם $f(x_1)=f(x_2)$, אבל $f$ יורדת ממש ולכן חד-חד-ערכית.
\end{solution}
